\documentclass[11pt]{book}
\usepackage{fontspec}
\usepackage{longtable,array,booktabs,xcolor,hyperref,xparse,geometry,microtype}
\geometry{a4paper,margin=1.75cm}
\setmainfont{FreeSerif}
\newfontfamily\devanagarifont[Script=Devanagari,Scale=1.05]{Noto Serif Devanagari}
\newcommand{\SA}[1]{{\devanagarifont #1}}
\hypersetup{unicode=true,colorlinks=true,linkcolor=blue,urlcolor=blue}
\setlength{\parindent}{0pt}
\setlength{\parskip}{0.4em}
\newcommand{\EmptyTarget}[1]{\hypertarget{#1}{}}
\newcommand{\XRef}[2]{\hyperlink{#1}{#2}}
\newcommand{\RuleRef}[1]{\XRef{rule:#1}{\texttt{#1}}}
\newcommand{\RootRef}[1]{\XRef{root:#1}{\texttt{#1}}}
\newcommand{\UseRef}[1]{\XRef{use:#1}{\texttt{#1}}}
\newcommand{\DerivRef}[1]{\XRef{deriv:#1}{\texttt{#1}}}
\newcommand{\BackToRule}[1]{\hfill{\scriptsize\XRef{rule:#1}{↑ global rule}}}
\newenvironment{bardown}[1]{\subsection*{#1}\begin{longtable}{>{\raggedleft\arraybackslash}p{0.29\linewidth}!{\vrule width .7pt}>{\raggedright\arraybackslash}p{0.66\linewidth}}}{\end{longtable}}
\newcommand{\BRow}[2]{#1 & #2\\[.7em]}
\newcommand{\LocalRuleUse}[4]{\EmptyTarget{use:#1}\textbf{\texttt{#1}} — \SA{#2}\par #3 \BackToRule{#4}}

\title{\SA{क्तप्रक्रियामानचित्रम्}\\A Process Map for Deciding the क्त Form of Any Sanskrit Dhātu}
\author{A generator-backed Bar-Down grammar archive}
\date{\today}
\begin{document}
\frontmatter
\maketitle
\tableofcontents

\chapter*{प्रयोगप्रतिज्ञा}
\addcontentsline{toc}{chapter}{प्रयोगप्रतिज्ञा}
This book is a derivational index. Explanations are stored once and reused by hyperlink. A chapter first links to a local “rules used” occurrence; that occurrence links to the ordered global rule registry. Root entries link back to the complete Dhātupāṭha registry.

\chapter*{\SA{निर्णयमानचित्रम्} — decision map}
\addcontentsline{toc}{chapter}{निर्णयमानचित्रम्}
\begin{bardown}{B0 — क्तनिर्णयस्य मुख्यप्रवाहः}
\BRow{\textbf{Input gate}}{Identify the exact Dhātupāṭha entry, not merely a dictionary spelling. Record gaṇa, indicatory markers, meaning, pada information, preverb, and intended syntactic interpretation.}
\BRow{\textbf{Affix gate}}{Invoke \RuleRef{AS-3-2-102}; register क्त as निष्ठा through \RuleRef{AS-1-1-26}.}
\BRow{\textbf{Voice gate}}{Test ordinary कर्मणि/resultative interpretation; then test \RuleRef{AS-3-4-72}.}
\BRow{\textbf{iṭ gate}}{Apply the ārdhadhātuka iṭ system beginning with \RuleRef{AS-7-2-35}; do not decide seṭ/aniṭ/veṭ from gaṇa alone.}
\BRow{\textbf{Stem gate}}{Apply root-specific and suffix-conditioned operations in grammatical order.}
\BRow{\textbf{Niṣṭhā gate}}{Search 8.2.42ff for substitution, optionality, prohibition, and meaning-conditioned forms.}
\BRow{\textbf{Sandhi gate}}{Apply only sandhi and tripādī operations actually triggered.}
\BRow{\textbf{Audit gate}}{Return surface form, derivation path, competing forms, interpretation, source status, and unresolved commentarial question.}
\end{bardown}


\mainmatter
\part{Preamble 1: \SA{धातुगणाः}}
\input{build/generated/dhatupatha.tex}

\part{Preamble 2: \SA{सूत्राणि}}
\input{build/generated/rules.tex}

\part{\SA{गणपदानि}}
\chapter{\SA{भ्वादिगणः}}\label{chapter:gana-01}
\section{\SA{सूत्राणि प्रयुक्तानि}}
\LocalRuleUse{G01-AS-1-1-26-U01}{क्तक्तवतू निष्ठा}{This local occurrence licenses the niṣṭhā designation.}{AS-1-1-26}
\LocalRuleUse{G01-AS-3-2-102-U01}{निष्ठा}{This local occurrence introduces क्त.}{AS-3-2-102}
\section{\SA{धातुपाठसम्बन्धः}}Return to \XRef{gana:01}{\SA{भ्वादिगणः}}.
\section{\SA{क्तनिर्माणबाराः}}\begin{bardown}{G01-B0}\BRow{Root identification}{Select a precise \RootRef{DHATU-01-XXXX} entry.}\BRow{Affix selection}{Use \UseRef{G01-AS-3-2-102-U01}.}\BRow{Niṣṭhā designation}{Use \UseRef{G01-AS-1-1-26-U01}.}\BRow{Expansion slot}{Insert audited root-specific bars here.}\end{bardown}

\chapter{\SA{अदादिगणः}}\label{chapter:gana-02}
\section{\SA{सूत्राणि प्रयुक्तानि}}\LocalRuleUse{G02-AS-1-1-26-U01}{क्तक्तवतू निष्ठा}{Local niṣṭhā designation.}{AS-1-1-26}\LocalRuleUse{G02-AS-3-2-102-U01}{निष्ठा}{Local क्त introduction.}{AS-3-2-102}
\section{\SA{धातुपाठसम्बन्धः}}Return to \XRef{gana:02}{\SA{अदादिगणः}}.
\section{\SA{क्तनिर्माणबाराः}}\begin{bardown}{G02-B0}\BRow{Root identification}{Select \RootRef{DHATU-02-XXXX}.}\BRow{Affix selection}{Use \UseRef{G02-AS-3-2-102-U01}.}\BRow{Expansion slot}{Insert audited bars.}\end{bardown}

\chapter{\SA{जुहोत्यादिगणः}}\label{chapter:gana-03}
\section{\SA{सूत्राणि प्रयुक्तानि}}\LocalRuleUse{G03-AS-1-1-26-U01}{क्तक्तवतू निष्ठा}{Local niṣṭhā designation.}{AS-1-1-26}\LocalRuleUse{G03-AS-3-2-102-U01}{निष्ठा}{Local क्त introduction.}{AS-3-2-102}
\section{\SA{धातुपाठसम्बन्धः}}Return to \XRef{gana:03}{\SA{जुहोत्यादिगणः}}.
\section{\SA{क्तनिर्माणबाराः}}\begin{bardown}{G03-B0}\BRow{Root identification}{Select \RootRef{DHATU-03-XXXX}.}\BRow{Affix selection}{Use \UseRef{G03-AS-3-2-102-U01}.}\BRow{Expansion slot}{Insert audited bars.}\end{bardown}

\chapter{\SA{दिवादिगणः}}\label{chapter:gana-04}
\section{\SA{सूत्राणि प्रयुक्तानि}}\LocalRuleUse{G04-AS-1-1-26-U01}{क्तक्तवतू निष्ठा}{Local niṣṭhā designation.}{AS-1-1-26}\LocalRuleUse{G04-AS-3-2-102-U01}{निष्ठा}{Local क्त introduction.}{AS-3-2-102}
\section{\SA{धातुपाठसम्बन्धः}}Return to \XRef{gana:04}{\SA{दिवादिगणः}}.
\section{\SA{क्तनिर्माणबाराः}}\begin{bardown}{G04-B0}\BRow{Root identification}{Select \RootRef{DHATU-04-XXXX}.}\BRow{Affix selection}{Use \UseRef{G04-AS-3-2-102-U01}.}\BRow{Expansion slot}{Insert audited bars.}\end{bardown}

\chapter{\SA{स्वादिगणः}}\label{chapter:gana-05}
\section{\SA{सूत्राणि प्रयुक्तानि}}\LocalRuleUse{G05-AS-1-1-26-U01}{क्तक्तवतू निष्ठा}{Local niṣṭhā designation.}{AS-1-1-26}\LocalRuleUse{G05-AS-3-2-102-U01}{निष्ठा}{Local क्त introduction.}{AS-3-2-102}
\section{\SA{धातुपाठसम्बन्धः}}Return to \XRef{gana:05}{\SA{स्वादिगणः}}.
\section{\SA{क्तनिर्माणबाराः}}\begin{bardown}{G05-B0}\BRow{Root identification}{Select \RootRef{DHATU-05-XXXX}.}\BRow{Affix selection}{Use \UseRef{G05-AS-3-2-102-U01}.}\BRow{Expansion slot}{Insert audited bars.}\end{bardown}

\chapter{\SA{तुदादिगणः}}\label{chapter:gana-06}
\section{\SA{सूत्राणि प्रयुक्तानि}}\LocalRuleUse{G06-AS-1-1-26-U01}{क्तक्तवतू निष्ठा}{Local niṣṭhā designation.}{AS-1-1-26}\LocalRuleUse{G06-AS-3-2-102-U01}{निष्ठा}{Local क्त introduction.}{AS-3-2-102}
\section{\SA{धातुपाठसम्बन्धः}}Return to \XRef{gana:06}{\SA{तुदादिगणः}}.
\section{\SA{क्तनिर्माणबाराः}}\begin{bardown}{G06-B0}\BRow{Root identification}{Select \RootRef{DHATU-06-XXXX}.}\BRow{Affix selection}{Use \UseRef{G06-AS-3-2-102-U01}.}\BRow{Expansion slot}{Insert audited bars.}\end{bardown}

\chapter{\SA{रुधादिगणः}}\label{chapter:gana-07}
\section{\SA{सूत्राणि प्रयुक्तानि}}\LocalRuleUse{G07-AS-1-1-26-U01}{क्तक्तवतू निष्ठा}{Local niṣṭhā designation.}{AS-1-1-26}\LocalRuleUse{G07-AS-3-2-102-U01}{निष्ठा}{Local क्त introduction.}{AS-3-2-102}
\section{\SA{धातुपाठसम्बन्धः}}Return to \XRef{gana:07}{\SA{रुधादिगणः}}.
\section{\SA{क्तनिर्माणबाराः}}\begin{bardown}{G07-B0}\BRow{Root identification}{Select \RootRef{DHATU-07-XXXX}.}\BRow{Affix selection}{Use \UseRef{G07-AS-3-2-102-U01}.}\BRow{Expansion slot}{Insert audited bars.}\end{bardown}

\chapter{\SA{तनादिगणः}}\label{chapter:gana-08}
\section{\SA{सूत्राणि प्रयुक्तानि}}\LocalRuleUse{G08-AS-1-1-26-U01}{क्तक्तवतू निष्ठा}{Local niṣṭhā designation.}{AS-1-1-26}\LocalRuleUse{G08-AS-3-2-102-U01}{निष्ठा}{Local क्त introduction.}{AS-3-2-102}
\section{\SA{धातुपाठसम्बन्धः}}Return to \XRef{gana:08}{\SA{तनादिगणः}}.
\section{\SA{क्तनिर्माणबाराः}}\begin{bardown}{G08-B0}\BRow{Root identification}{Select \RootRef{DHATU-08-XXXX}.}\BRow{Affix selection}{Use \UseRef{G08-AS-3-2-102-U01}.}\BRow{Expansion slot}{Insert audited bars.}\end{bardown}

\chapter{\SA{क्र्यादिगणः}}\label{chapter:gana-09}
\section{\SA{सूत्राणि प्रयुक्तानि}}\LocalRuleUse{G09-AS-1-1-26-U01}{क्तक्तवतू निष्ठा}{Local niṣṭhā designation.}{AS-1-1-26}\LocalRuleUse{G09-AS-3-2-102-U01}{निष्ठा}{Local क्त introduction.}{AS-3-2-102}
\section{\SA{धातुपाठसम्बन्धः}}Return to \XRef{gana:09}{\SA{क्र्यादिगणः}}.
\section{\SA{क्तनिर्माणबाराः}}\begin{bardown}{G09-B0}\BRow{Root identification}{Select \RootRef{DHATU-09-XXXX}.}\BRow{Affix selection}{Use \UseRef{G09-AS-3-2-102-U01}.}\BRow{Expansion slot}{Insert audited bars.}\end{bardown}

\chapter{\SA{चुरादिगणः}}\label{chapter:gana-10}
\section{\SA{सूत्राणि प्रयुक्तानि}}\LocalRuleUse{G10-AS-1-1-26-U01}{क्तक्तवतू निष्ठा}{Local niṣṭhā designation.}{AS-1-1-26}\LocalRuleUse{G10-AS-3-2-102-U01}{निष्ठा}{Local क्त introduction.}{AS-3-2-102}
\section{\SA{धातुपाठसम्बन्धः}}Return to \XRef{gana:10}{\SA{चुरादिगणः}}.
\section{\SA{क्तनिर्माणबाराः}}\begin{bardown}{G10-B0}\BRow{Root identification}{Select \RootRef{DHATU-10-XXXX}.}\BRow{Affix selection}{Use \UseRef{G10-AS-3-2-102-U01}.}\BRow{Expansion slot}{Insert audited bars.}\end{bardown}

\chapter{\SA{कण्ड्वादिगणः}}\label{chapter:gana-11}
\section{\SA{सूत्राणि प्रयुक्तानि}}\LocalRuleUse{G11-AS-1-1-26-U01}{क्तक्तवतू निष्ठा}{Local niṣṭhā designation.}{AS-1-1-26}\LocalRuleUse{G11-AS-3-2-102-U01}{निष्ठा}{Local क्त introduction.}{AS-3-2-102}
\section{\SA{धातुपाठसम्बन्धः}}Return to \XRef{gana:11}{\SA{कण्ड्वादिगणः}}.
\section{\SA{क्तनिर्माणबाराः}}\begin{bardown}{G11-B0}\BRow{Root identification}{Select \RootRef{DHATU-11-XXXX}.}\BRow{Affix selection}{Use \UseRef{G11-AS-3-2-102-U01}.}\BRow{Expansion slot}{Insert audited bars.}\end{bardown}


\part{Cross-cutting registries}
\input{build/generated/categories.tex}
\chapter{\SA{सन्धिनियमाः प्रयुक्ताः}}
This registry contains only rules actually invoked in a chapter. New rules are added globally once, then cited through local usage anchors.
\section{\RuleRef{AS-8-4-55} — \SA{खरि च}}
Final devoicing and related tripādī behavior must be applied at the correct stage; it is not decorative spelling cleanup.

\input{build/generated/irregular.tex}

\backmatter
\chapter{Audit ledger}
\input{build/generated/audit.tex}
\end{document}
